\section{KDE Project Analysis} % {{{

\marginnote{\centering\includegraphics[width=0.5\marginparwidth]{kde/logo}}

KDE is a desktop environment designed to run on UNIX based systems as well as
Microsoft Windows and Mac \textsmaller{OS X} systems \cite{KDEPress,KDEAbout}.
Since 2010 and the release of KDE~4.4 the project is known as \ac{KDE SC}. It
is composed by the desktop environment named Plasma Desktop and several core
applications for the daily needs. \ac{KDE SC} is licensed under the \ac{GPL}
and \ac{LGPL}. The name was originally an acronym for Kool Desktop Environment
and later for K Desktop Environment, however it is no longer in use. KDE is a
modular project consisting of several libraries which allow developers to build
their programs around the platform. As the KDE project and the \ac{KDE SC}
stand for a wide range of technologies and applications, this analysis will be
limited to the core of the \ac{KDE SC}, also known as KDE Base.

\subsection{History} % {{{

KDE was first announced by Matthias Ettrich in 1996 as a desktop environment
for end users with the same look and feel for all applications
\cite{KDEAnnouncement}. The name KDE was a word play with the existing
\ac{CDE}, which was highly popular at that time. Matthias Ettrich chose the Qt
framework as the graphical library, which was and is developed by the company
Trolltech. KDE~1.0 was then released in 1998 \cite{KDEHistory}. At the same
time, Trolltech dual licensed the framework under the \ac{QPL} and a
proprietary software license. However, it was debated if that license was
compatible with the \ac{GPL} and in 2000 the framework was licensed under the
\ac{GPL} which ceased the criticism.

In 2009 the KDE project renamed the project to \ac{KDE SC} and redefined the
project as a community which delivers free software for user interfaces by
emphasizing the KDE technologies \cite{KDESC}. All software created with KDE
technologies are KDE projects. However \ac{KDE SC} only contains projects which
derive from the project itself and share a common release cycle.

\ac{KDE SC}~4 was a new approach to the desktop metaphor and created a new user
experience for users. The centerpiece is the Plasma Workspace which exists for
several devices, such as for desktop computers, netbooks, tablets and
smartphones.

% }}}

\subsection{Community} % {{{

The KDE project is one of the largest \ac{FOSS} projects and according to the
community the second largest after the Linux Kernel \cite{KDEPress}. There are
about 1800 active contributors to the project and its surroundings and is used
by over a million people. Also, the community spans around projects based on
KDE technologies.

Most of the communication in the KDE project takes place via mailing lists,
most importantly \emph{kde-devel} and \emph{kde-core-devel}
\cite{KDEProjectManagement,KDEContribute}. While the first one is mostly used
for communication by application developers, the latter one is for
communication on the KDE core project. Furthermore communication can happen
over \ac{IRC} channels, blogs or forums.

The most important KDE conference is the annual \emph{Akademy} which is held at
different locations in Europe \cite{KDEHistory}. The first KDE conferences were
however named after the current release which started with \emph{KDE One} in
1997. That is also the reason, why the conferences did not take place annually
at that time. In 2004 it changed the name to Akademy. Beginning with 2009 the
project incorporated the conference into the \emph{Desktop Summit} in
collaboration with the GNOME project.

\begin{table}
  \centering
  \begin{tabularx}{\textwidth}{lXr}
    \toprule
    \tableheadline{Event}   & \tableheadline{Venue}           & \tableheadline{Date} \\
    \midrule
    KDE One                 & Arnsberg, Germany               & 1997 \\
    KDE Two                 & Erlangen, Germany               & 1999 \\
    KDE Three Beta          & Trysil, Norway                  & 2000 \\
    KDE Three               & Nürnberg, Germany               & 2002 \\
    Kastle                  & Nové Hrady, Czech Republic      & 2003 \\
    aKademy                 & Ludwigsburg, Germany            & 2004 \\
    aKademy                 & Málaga, Spain                   & 2005 \\
    aKademy                 & Dublin, Ireland                 & 2006 \\
    aKademy                 & Glasgow, Scotland               & 2007 \\
    Akademy                 & Sint-Katelijne-Waver, Belgium   & 2008 \\
    Desktop Summit          & Gran Canaria, Spain             & 2009 \\
    Akademy                 & Tampere, Finland                & 2010 \\
    Desktop Summit          & Berlin, Germany                 & 2011 \\
    Akademy                 & Tallinn, Estonia                & 2012 \\
    \bottomrule
  \end{tabularx}
  \caption[Previous and Planned KDE Conferences]{Previous and planned KDE conferences.}
\end{table}

\begin{figure}[htbp]
  \centering
  \includegraphics[width=0.8\textwidth]{kde/commits_by_author}
  \caption[Commits by Most Active Authors, KDE]
  {Monthly activity of the most active KDE Base developers. The most active
    developers almost switched completely when comparing the inception phase to
    nowadays.}
  \label{fig:kde:cba}
\end{figure}

Due to the size of the KDE project it is organized around many independent
teams with a communication structure between them. There are few groups who
stand above those teams and only take care about coordinating the project
\cite{KDEDevelopmentModel,KDEProjectManagement}.

\paragraph{Contributor}

After having contributed to any KDE project for some time and planning to
contribute in the future, a KDE contributor account can be requested
\cite{KDEContribute,KDEContributor}. In the request one must state why one is
interested in having an account. The KDE sysadmin team will then check the
application and grant an account or not.

\paragraph{Module Coordinator}

As the KDE project is highly modular, each application or library has an
independent team with its own structure besides a few exceptions. A module
coordinator plans and executes releases and the direction of a single module.
As the teams are so diverse, there are many ways to become a module
coordinator. The most obvious is of course by unweary contributions to a
project.

\begin{figure}[htbp]
  \centering
  \includegraphics[width=0.8\textwidth]{kde/commits_by_year}
  \caption[Commits by Year, KDE]
  {Yearly overview of commits to KDE Base. The development and release of KDE~4
    and following releases in 2008 certainly attracted a lot of attention.}
  \label{fig:kde:cby}
\end{figure}

\paragraph{Core Team}

One of the most important teams is the core team as it defines the overall
direction the project is heading \cite{KDEProjectManagement}. There is no
single person inside the team responsible for decision making, instead the team
discusses the issue and comes up with a solution together. The primary
communication method is the \emph{kde-core-devel} mailing list which is
publicly readable, however an approval is required to join it. The KDE project
does not define clear restrictions to membership and therefore a membership is
granted by invite only after distinguishing or outstanding work for the KDE
project.

\paragraph{Release Team}

The actual release and the release schedules are provided and enforced by the
release team \cite{KDEReleaseTeam}. It is composed by module coordinators and
other release team members who actually implement the releases. The release
team makes sure that all modules are following the release schedule and are in
a good shape. Also they decide on code freeze breaks and take decision about
future features of a release.

\paragraph{KDE Founder}

Matthias Ettrich, the founder of the KDE project is no longer actively involved
in the project as he works full time on the underlying Qt framework. However he
has still an important voice in the community and contributes indirectly to KDE
by his role in the Qt project.

% }}}

\subsection{Release Process} % {{{

\begin{figure}[htbp]
  \centering
  \includegraphics[width=\textwidth]{kde/releases}
  \caption[Major Releases of KDE]{Platform and standard releases of KDE.}
\end{figure}

The KDE project uses a versioning scheme with three numbers and distinguishes
between major and minor releases
\cite{KDEReleaseTeam,KDEReleaseSchedule,KDESchedule}. The numbers are called
major, minor and micro numbers. The major and minor numbers define major
releases while the micro number defines maintenance releases. The release is a
platform or standard release depending on whether the major or minor numbers
change. A platform release defines a new series of releases which can break
backwards compatibility to the previous platform release. They are often
planned a long time ahead, usually including changes of used libraries or used
versions switches. A standard release however has to maintain backwards
compatibility with its series. They can have new features and user interface
changes, however a KDE application has to work on all releases of the same
platform. Maintenance releases are not allowed to come with new features,
except bug fixes or small enhancements. Beginning with \ac{KDE SC}~4 the cycle
changed to a major release every six months and a maintenance release roughly
every month except when a major release comes out.

To ensure the release schedule, the release team publishes a new release
schedule after each major release. The schedule includes all deadlines and
freezes which will lead to the next major release and includes future
maintenance releases. This means of course that two release schedules will
overlap during the month before the new major release. The following freezes
and releases are provided and enforced by the release team
\cite{KDEReleaseSchedule}.

\paragraph{Soft Feature Freeze}

Approximately three weeks before the first beta release no new features are
allowed except already approved and planned features. Not finalized features
have to be postponed to the next major release.

\begin{figure}[hbtp]
  \centering
  \includegraphics[width=\textwidth]{kde/punchcard}
  \caption[Time Based View on Commits, KDE]
  {Time based view on commits of core contributors. The times are almost
    equally distributed probably meaning that there exist a lot of employed and
    volunteering developers in the community.}
  \label{fig:kde:p}
\end{figure}

\begin{figure}[htbp]
  \centering
  \includegraphics[width=0.8\textwidth]{kde/commits_by_month}
  \caption[Commits by Month, KDE]
  {Amount of commits per month of core contributors. Again, the development
    phase of KDE~4 and following releases is quite visible.}
  \label{fig:kde:cbm}
\end{figure}

\paragraph{Dependency Freeze}

Approximately two weeks before the first beta release no new or additional
dependency versions are allowed anymore. However it is possible to request an
exception by the release team.

\paragraph{Soft Message, Soft API and Hard Feature Freeze}

One week before the first beta release no new strings and changes to the
existing messages can be made except corrections. Additionally the existing
\ac{API} should be almost finished and if changes are made they should be
reported to the relevant teams. Lastly no new features can be added to the
projects, even if they were planned for this release cycle.

\paragraph{Beta Release}

Approximately two months before the major release the first beta release gets
published. This is usually followed by a second beta release two weeks later.

\paragraph{Hard API, Message and Documentation Freeze}

Five weeks before the major release the \ac{API}, translatable messages and the
documentation must be finished and can't be changed anymore.

\paragraph{Release Candidate}

Approximately a month before the final release the first release candidate will
be published. Two weeks later the second release candidate follows.

\begin{figure}[htbp]
  \centering
  \includegraphics[width=0.8\textwidth]{kde/authors_by_month}
  \caption[Authors by Month, KDE]
  {Amount of distinct authors of KDE Base over time. The project appears to
    grow quite linearly with a small down before the start of the KDE~4
    development and the appropriate boost during and after the development.}
  \label{fig:kde:abm}
\end{figure}

\paragraph{Final Release}

The final release will then be published three weeks after the last release
candidate.

\paragraph{Minor Releases}

A new minor release will be published each following month until the next major
release gets published. This normally leads to four minor releases.

% }}}

\subsection{Development} % {{{

During the development of a new major release, new project goals are set using
feature proposals \cite{KDEDevelopmentModel,KDEFAQ}. This normally happens
after each major release when KDE developers have the time to propose new
features for the next major release. Each feature must have an author who is
willing to implement it. This often limits the number of features to already
contributing KDE developers. A KDE developer is free to implement the feature
by himself, however the process is often accompanied by a discussion with the
specific team and several adjustments to the implementation.

During the soft feature freeze is decided whether the feature will be available
for the upcoming major release or if it will be postponed to the next major
release. If the implementation already started and is almost finished, it can
continue. Otherwise it will be postponed. This development process is a quite
deliberate approach where people are free to implement what they see missing in
the project.

% }}}

% }}}
